% ============================================
% DOC-23 (ZH)
% 中文翻译 — SPEC-935-R356
% ============================================
\section*{\textbf{DOC-23} --- 档案馆事件记录}
\markboth{DOC-23 --- 档案馆事件记录}{}

\noindent \textit{从鉴定档案中提取的文件。}

\vspace{0.5em}\noindent\rule{\textwidth}{0.5pt}\vspace{0.5em}


\noindent {\footnotesize\texttt{公共档案馆 $\cdot$ 书证 $\cdot$ 案卷1,260}}\\[-0.2em]

\noindent \textbf{米纳斯吉拉斯州公共档案馆 —— 内部安保报告}\\
\noindent \textbf{记录号:} 2026-RO-0847\\
\noindent\textbf{日期:} 2026年8月3日。 \textbf{时间:} 8点15分。

\noindent \textbf{申请单位:} 特藏部管理科。\\
\noindent \textbf{地点:} 保险柜室 —— 地下，C区。

\noindent \textbf{事件经过:}

\noindent 8月3日7点50分，值夜班的警卫——埃德松·莫雷拉——向值班中心报告，永久保险柜的重量传感器在凌晨曾触发。传感器检测内容物总质量的变化。警报于3点14分响起，持续六秒，然后自行静默。值班中心没有派遣警卫前往现场，因为重量传感器警报规程要求目视确认，但保险柜内部摄像头在触发期间没有显示与基线不同的画面。

\noindent 8点整，馆藏经理与安保技术员打开保险柜进行查验。

\noindent \textbf{查明:}

\noindent a) 1,260物品在原有位置，保存包装内；

\noindent b) 包装的安全封条完好，与封存时登记编号一致（封条号 APM-2026-4471）；

\noindent c) 室内精密天平显示，该物品\textbf{比2026年7月28日入库时重了12克}。设备误差范围为0.2克。

\noindent \textbf{已采取措施:}

\noindent a) 用独立天平复称 —— 已确认：增加12克；

\noindent b) 重新封存包装，并对重量和旧封条进行拍照记录；

\noindent c) 通知管理局，管理局决定对事件保持行政沉默；

\noindent d) 应其本人请求更换夜班警卫。埃德松·莫雷拉在离职时声明：“不是天平的问题。是包装里面的东西动了。我听见了。纸自己动不会发出声音。”

\noindent 在签署离职文件之前，埃德松补充道——安保技术员出于谨慎记录了这段话：\textquotedblleft{}我家里有两个孩子和一个妻子。无论给多少工资，我夜里都不会再回那个房间。你们可以说那是潮气，是传感器故障，随便说什么。我决定宁可不知道那是什么。我宁可睡个安稳觉。\textquotedblright{}他的恐惧在标准表格上没有任何一栏可填，但档案主管在归档副本的页边手写道：\textquotedblleft{}他没有说谎。\textquotedblright{}

\noindent \textbf{归档于:} 2026年8月3日，按管理局决定。受限访问。





\vspace{0.5em}
\noindent\textit{\textbf{↪ 路线:} \rota{DOC-01} \rotasep \rota{LUC-02} \rotasep \rota{DOC-24}}
